\documentclass[aspectratio=169,11pt]{beamer}
\usepackage[T1]{fontenc}
\usepackage{sourceserifpro}
\usepackage[default]{sourcesanspro}
\usepackage{booktabs}
\usepackage{graphicx}
\usepackage{xcolor}
\usepackage{tikz}

% Generated by build_numbers.py from results/*.json. Do not edit.
\newcommand{\AdvIllicitAUC}{0.998}
\newcommand{\AdvIllicitClockAUC}{1.000}
\newcommand{\AdvIllicitClockKSShare}{73\%}
\newcommand{\AdvIllicitKSShare}{72\%}
\newcommand{\AdvIllicitN}{169}
\newcommand{\AdvLicitAUC}{0.998}
\newcommand{\AdvLicitClockAUC}{1.000}
\newcommand{\AdvLicitClockKSShare}{10\%}
\newcommand{\AdvLicitKSShare}{7\%}
\newcommand{\AdvLicitN}{6,518}
\newcommand{\ClockAblRFFone}{0.801}
\newcommand{\ClockAblRFFoneDiff}{0.000}
\newcommand{\ClockAblRFFoneP}{=1.00}
\newcommand{\ClockAblRFPRAUCDiff}{0.000}
\newcommand{\ClockAblRFPRAUCP}{=0.84}
\newcommand{\ClockAblXGBFone}{0.800}
\newcommand{\ClockAblXGBFoneDiff}{-0.000}
\newcommand{\ClockAblXGBFoneP}{=0.85}
\newcommand{\ClockAblXGBPRAUCDiff}{0.001}
\newcommand{\ClockAblXGBPRAUCP}{=0.26}
\newcommand{\ClockCount}{5}
\newcommand{\ClockList}{agg\_07, agg\_09, agg\_42, agg\_06, agg\_45}
\newcommand{\ClockMinRho}{0.953}
\newcommand{\DriftDirSAGEPost}{0.016}
\newcommand{\DriftDirSAGEPre}{0.687}
\newcommand{\DriftLogRegPost}{0.066}
\newcommand{\DriftLogRegPre}{0.582}
\newcommand{\DriftMLPPost}{0.031}
\newcommand{\DriftMLPPre}{0.720}
\newcommand{\DriftRFPost}{0.043}
\newcommand{\DriftRFPre}{0.890}
\newcommand{\DriftXGBPost}{0.036}
\newcommand{\DriftXGBPre}{0.888}
\newcommand{\EdaComponents}{49}
\newcommand{\EdaFeaturesAfterFilter}{103}
\newcommand{\EdaHomBase}{9.8\%}
\newcommand{\EdaHomIll}{54.1\%}
\newcommand{\EdaHomLic}{2.4\%}
\newcommand{\EdaIllicitPost}{169}
\newcommand{\EdaIllicitPre}{914}
\newcommand{\EdaIllicitShare}{9.8\%}
\newcommand{\EdaPairsNinetyFive}{96}
\newcommand{\EdaShareTest}{6.5\%}
\newcommand{\EdaShareTrain}{11.6\%}
\newcommand{\GraphDirSAGEgraphFone}{0.669}
\newcommand{\GraphDirSAGEgraphFoneCI}{[0.439,\, 0.757]}
\newcommand{\GraphDirSAGEgraphPRAUC}{0.702}
\newcommand{\GraphGCNgraphFone}{0.515}
\newcommand{\GraphGCNgraphFoneCI}{[0.274,\, 0.644]}
\newcommand{\GraphGCNgraphPRAUC}{0.608}
\newcommand{\GraphPairedRFAFGFoneDiff}{0.014}
\newcommand{\GraphPairedRFAFGFoneP}{<0.001}
\newcommand{\GraphPairedRFAFGNEFoneDiff}{0.007}
\newcommand{\GraphPairedRFAFGNEFoneP}{=0.14}
\newcommand{\GraphPairedRFAFGNEPRAUCDiff}{-0.004}
\newcommand{\GraphPairedRFAFGNEPRAUCP}{=0.003}
\newcommand{\GraphPairedRFAFGPRAUCDiff}{-0.000}
\newcommand{\GraphPairedRFAFGPRAUCP}{=0.19}
\newcommand{\GraphPairedRFAFNEFoneDiff}{0.013}
\newcommand{\GraphPairedRFAFNEFoneP}{<0.001}
\newcommand{\GraphPairedRFAFNEPRAUCDiff}{-0.004}
\newcommand{\GraphPairedRFAFNEPRAUCP}{=0.01}
\newcommand{\GraphPairedRFLFFoneDiff}{-0.061}
\newcommand{\GraphPairedRFLFFoneP}{<0.001}
\newcommand{\GraphPairedRFLFPRAUCDiff}{-0.009}
\newcommand{\GraphPairedRFLFPRAUCP}{<0.001}
\newcommand{\GraphPairedXGBAFGFoneDiff}{-0.009}
\newcommand{\GraphPairedXGBAFGFoneP}{=0.001}
\newcommand{\GraphPairedXGBAFGNEFoneDiff}{0.017}
\newcommand{\GraphPairedXGBAFGNEFoneP}{=0.04}
\newcommand{\GraphPairedXGBAFGNEPRAUCDiff}{-0.003}
\newcommand{\GraphPairedXGBAFGNEPRAUCP}{=0.36}
\newcommand{\GraphPairedXGBAFGPRAUCDiff}{0.002}
\newcommand{\GraphPairedXGBAFGPRAUCP}{=0.04}
\newcommand{\GraphPairedXGBAFNEFoneDiff}{0.014}
\newcommand{\GraphPairedXGBAFNEFoneP}{=0.06}
\newcommand{\GraphPairedXGBAFNEPRAUCDiff}{-0.001}
\newcommand{\GraphPairedXGBAFNEPRAUCP}{=0.77}
\newcommand{\GraphPairedXGBLFFoneDiff}{-0.047}
\newcommand{\GraphPairedXGBLFFoneP}{<0.001}
\newcommand{\GraphPairedXGBLFPRAUCDiff}{-0.015}
\newcommand{\GraphPairedXGBLFPRAUCP}{=0.001}
\newcommand{\GraphRFAFFone}{0.804}
\newcommand{\GraphRFAFFoneCI}{[0.618,\, 0.888]}
\newcommand{\GraphRFAFGFone}{0.813}
\newcommand{\GraphRFAFGFoneCI}{[0.639,\, 0.896]}
\newcommand{\GraphRFAFGNEFone}{0.808}
\newcommand{\GraphRFAFGNEFoneCI}{[0.627,\, 0.890]}
\newcommand{\GraphRFAFGNEPRAUC}{0.789}
\newcommand{\GraphRFAFGPRAUC}{0.793}
\newcommand{\GraphRFAFNEFone}{0.814}
\newcommand{\GraphRFAFNEFoneCI}{[0.636,\, 0.894]}
\newcommand{\GraphRFAFNEPRAUC}{0.789}
\newcommand{\GraphRFAFPRAUC}{0.793}
\newcommand{\GraphRFLFFone}{0.743}
\newcommand{\GraphRFLFFoneCI}{[0.548,\, 0.839]}
\newcommand{\GraphRFLFGFone}{0.748}
\newcommand{\GraphRFLFGFoneCI}{[0.552,\, 0.845]}
\newcommand{\GraphRFLFGPRAUC}{0.785}
\newcommand{\GraphRFLFPRAUC}{0.784}
\newcommand{\GraphSkipGCNgraphFone}{0.507}
\newcommand{\GraphSkipGCNgraphFoneCI}{[0.261,\, 0.639]}
\newcommand{\GraphSkipGCNgraphPRAUC}{0.624}
\newcommand{\GraphXGBAFFone}{0.794}
\newcommand{\GraphXGBAFFoneCI}{[0.613,\, 0.886]}
\newcommand{\GraphXGBAFGFone}{0.787}
\newcommand{\GraphXGBAFGFoneCI}{[0.604,\, 0.878]}
\newcommand{\GraphXGBAFGNEFone}{0.816}
\newcommand{\GraphXGBAFGNEFoneCI}{[0.628,\, 0.896]}
\newcommand{\GraphXGBAFGNEPRAUC}{0.795}
\newcommand{\GraphXGBAFGPRAUC}{0.800}
\newcommand{\GraphXGBAFNEFone}{0.812}
\newcommand{\GraphXGBAFNEFoneCI}{[0.626,\, 0.894]}
\newcommand{\GraphXGBAFNEPRAUC}{0.797}
\newcommand{\GraphXGBAFPRAUC}{0.798}
\newcommand{\GraphXGBLFFone}{0.750}
\newcommand{\GraphXGBLFFoneCI}{[0.568,\, 0.840]}
\newcommand{\GraphXGBLFGFone}{0.749}
\newcommand{\GraphXGBLFGFoneCI}{[0.560,\, 0.841]}
\newcommand{\GraphXGBLFGPRAUC}{0.789}
\newcommand{\GraphXGBLFPRAUC}{0.782}
\newcommand{\HybRFFone}{0.664}
\newcommand{\HybRFFoneDiff}{-0.127}
\newcommand{\HybRFFoneDrop}{0.127}
\newcommand{\HybRFFoneP}{<0.001}
\newcommand{\HybRFPRAUC}{0.696}
\newcommand{\HybXGBFone}{0.662}
\newcommand{\HybXGBFoneDiff}{-0.117}
\newcommand{\HybXGBFoneDrop}{0.117}
\newcommand{\HybXGBFoneP}{=0.001}
\newcommand{\HybXGBPRAUC}{0.692}
\newcommand{\IFAFLicitAnomalousPRAUC}{0.036}
\newcommand{\IFAFLicitAnomalousROCAUC}{0.176}
\newcommand{\IFAFLicitTypicalPRAUC}{0.473}
\newcommand{\IFAFLicitTypicalROCAUC}{0.824}
\newcommand{\IFAFUnsupAnomalousPRAUC}{0.036}
\newcommand{\IFAFUnsupAnomalousROCAUC}{0.170}
\newcommand{\IFAFUnsupTypicalPRAUC}{0.475}
\newcommand{\IFAFUnsupTypicalROCAUC}{0.830}
\newcommand{\IFLFLicitAnomalousPRAUC}{0.037}
\newcommand{\IFLFLicitAnomalousROCAUC}{0.167}
\newcommand{\IFLFLicitTypicalPRAUC}{0.616}
\newcommand{\IFLFLicitTypicalROCAUC}{0.833}
\newcommand{\IFLFUnsupAnomalousPRAUC}{0.038}
\newcommand{\IFLFUnsupAnomalousROCAUC}{0.181}
\newcommand{\IFLFUnsupTypicalPRAUC}{0.374}
\newcommand{\IFLFUnsupTypicalROCAUC}{0.819}
\newcommand{\ImbLRNoneFixedFone}{0.461}
\newcommand{\ImbLRNonePRAUC}{0.319}
\newcommand{\ImbLRNoneTunedFone}{0.515}
\newcommand{\ImbLRSmoteFixedFone}{0.350}
\newcommand{\ImbLRSmotePRAUC}{0.275}
\newcommand{\ImbLRSmoteTunedFone}{0.470}
\newcommand{\ImbLRWeightFixedFone}{0.329}
\newcommand{\ImbLRWeightPRAUC}{0.270}
\newcommand{\ImbLRWeightTunedFone}{0.464}
\newcommand{\ImbMLPNoneFixedFone}{0.620}
\newcommand{\ImbMLPNonePRAUC}{0.605}
\newcommand{\ImbMLPNoneTunedFone}{0.679}
\newcommand{\ImbMLPSmoteFixedFone}{0.620}
\newcommand{\ImbMLPSmotePRAUC}{0.601}
\newcommand{\ImbMLPSmoteTunedFone}{0.651}
\newcommand{\ImbMLPWeightFixedFone}{0.621}
\newcommand{\ImbMLPWeightPRAUC}{0.587}
\newcommand{\ImbMLPWeightTunedFone}{0.640}
\newcommand{\ImbRFNoneFixedFone}{0.827}
\newcommand{\ImbRFNonePRAUC}{0.793}
\newcommand{\ImbRFNoneTunedFone}{0.804}
\newcommand{\ImbRFSmoteFixedFone}{0.798}
\newcommand{\ImbRFSmotePRAUC}{0.807}
\newcommand{\ImbRFSmoteTunedFone}{0.775}
\newcommand{\ImbRFWeightFixedFone}{0.812}
\newcommand{\ImbRFWeightPRAUC}{0.796}
\newcommand{\ImbRFWeightTunedFone}{0.788}
\newcommand{\ImbXGBNoneFixedFone}{0.818}
\newcommand{\ImbXGBNonePRAUC}{0.798}
\newcommand{\ImbXGBNoneTunedFone}{0.794}
\newcommand{\ImbXGBSmoteFixedFone}{0.818}
\newcommand{\ImbXGBSmotePRAUC}{0.816}
\newcommand{\ImbXGBSmoteTunedFone}{0.808}
\newcommand{\ImbXGBWeightFixedFone}{0.805}
\newcommand{\ImbXGBWeightPRAUC}{0.800}
\newcommand{\ImbXGBWeightTunedFone}{0.805}
\newcommand{\LabelEffA}{0.361}
\newcommand{\LabelEffB}{0.385}
\newcommand{\LabelEffC}{0.542}
\newcommand{\LabelEffD}{0.703}
\newcommand{\LabelEffE}{0.714}
\newcommand{\LabelEffF}{0.745}
\newcommand{\LabelEffG}{0.761}
\newcommand{\LabelEffH}{0.768}
\newcommand{\LabelEffI}{0.785}
\newcommand{\LabelEffJ}{0.794}
\newcommand{\LabelEffK}{0.798}
\newcommand{\LabelEffNA}{15}
\newcommand{\LabelEffNB}{30}
\newcommand{\LabelEffNC}{60}
\newcommand{\LabelEffND}{149}
\newcommand{\LabelEffNE}{299}
\newcommand{\LabelEffNF}{598}
\newcommand{\LabelEffNG}{1,495}
\newcommand{\LabelEffNH}{2,989}
\newcommand{\LabelEffNI}{7,474}
\newcommand{\LabelEffNJ}{14,947}
\newcommand{\LabelEffNK}{29,894}
\newcommand{\LabelsFullPRAUC}{0.798}
\newcommand{\LabelsToBeatIF}{60}
\newcommand{\MainLRBudget}{39\%}
\newcommand{\MainLRFone}{0.515}
\newcommand{\MainLRFoneCI}{[0.328,\, 0.626]}
\newcommand{\MainLRPRAUC}{0.319}
\newcommand{\MainLRPRAUCCI}{[0.173,\, 0.443]}
\newcommand{\MainLRPrecision}{0.417}
\newcommand{\MainLRRecall}{0.674}
\newcommand{\MainMLPBudget}{57\%}
\newcommand{\MainMLPFone}{0.679}
\newcommand{\MainMLPFoneCI}{[0.481,\, 0.750]}
\newcommand{\MainMLPPRAUC}{0.605}
\newcommand{\MainMLPPRAUCCI}{[0.425,\, 0.763]}
\newcommand{\MainMLPPrecision}{0.739}
\newcommand{\MainMLPRecall}{0.632}
\newcommand{\MainRFBudget}{61\%}
\newcommand{\MainRFFone}{0.804}
\newcommand{\MainRFFoneCI}{[0.618,\, 0.888]}
\newcommand{\MainRFPRAUC}{0.793}
\newcommand{\MainRFPRAUCCI}{[0.576,\, 0.900]}
\newcommand{\MainRFPrecision}{0.896}
\newcommand{\MainRFRecall}{0.730}
\newcommand{\MainXGBBudget}{62\%}
\newcommand{\MainXGBFone}{0.794}
\newcommand{\MainXGBFoneCI}{[0.613,\, 0.886]}
\newcommand{\MainXGBPRAUC}{0.798}
\newcommand{\MainXGBPRAUCCI}{[0.586,\, 0.903]}
\newcommand{\MainXGBPrecision}{0.865}
\newcommand{\MainXGBRecall}{0.734}
\newcommand{\MemDirSAGERFImp}{99.8\%}
\newcommand{\MemDirSAGETestPRAUC}{0.704}
\newcommand{\MemDirSAGETrainPRAUC}{0.9999}
\newcommand{\MemGCNRFImp}{14.0\%}
\newcommand{\MemGCNTestPRAUC}{0.584}
\newcommand{\MemGCNTrainPRAUC}{0.9797}
\newcommand{\NovelPost}{90\%}
\newcommand{\NovelPre}{31\%}
\newcommand{\PairedXGBLRFoneDiff}{0.284}
\newcommand{\PairedXGBLRFoneP}{<0.001}
\newcommand{\PairedXGBLRPRAUCP}{<0.001}
\newcommand{\PairedXGBMLPFoneDiff}{0.133}
\newcommand{\PairedXGBMLPFoneP}{<0.001}
\newcommand{\PairedXGBMLPPRAUCP}{<0.001}
\newcommand{\PairedXGBRFFoneDiff}{-0.003}
\newcommand{\PairedXGBRFFoneP}{=0.50}
\newcommand{\PairedXGBRFPRAUCP}{=0.07}
\newcommand{\PermAggregated}{0.059}
\newcommand{\PermLocal}{0.695}
\newcommand{\PermStructure}{0.001}
\newcommand{\RepLRAF}{0.449}
\newcommand{\RepLRLF}{0.361}
\newcommand{\RepMLPAF}{0.577}
\newcommand{\RepMLPAFSd}{0.026}
\newcommand{\RepMLPLF}{0.643}
\newcommand{\RepMLPLFSd}{0.036}
\newcommand{\RepRFAF}{0.810}
\newcommand{\RepRFAFSd}{0.005}
\newcommand{\RepRFLF}{0.778}
\newcommand{\RepRFLFSd}{0.008}
\newcommand{\ShapAgg}{35\%}
\newcommand{\ShapLocal}{62\%}
\newcommand{\ShapStruct}{3\%}
\newcommand{\ShapTopFeature}{local\_52}
\newcommand{\UnsBaseRate}{0.065}
\newcommand{\WFExpandingPost}{0.185}
\newcommand{\WFExpandingPre}{0.905}
\newcommand{\WFSlidingPost}{0.144}
\newcommand{\WFSlidingPre}{0.896}
\newcommand{\WFStaticPost}{0.043}
\newcommand{\WFStaticPre}{0.890}


% ---------------------------------------------------------------- theme (matches the pitch)
\definecolor{paper}{HTML}{FAFAF7}
\definecolor{cover}{HTML}{F0EFEA}
\definecolor{ink}{HTML}{222222}
\definecolor{muted}{HTML}{6B6B6B}
\definecolor{accent}{HTML}{1F4E8C}
\definecolor{illicit}{HTML}{B03A2E}
\definecolor{rule}{HTML}{DADADA}

\setbeamercolor{background canvas}{bg=paper}
\setbeamercolor{normal text}{fg=ink}
\setbeamercolor{frametitle}{fg=ink}
\setbeamercolor{itemize item}{fg=ink}
\setbeamercolor{enumerate item}{fg=ink}
\setbeamercolor{alerted text}{fg=illicit}
\setbeamertemplate{navigation symbols}{}
\setbeamertemplate{itemize items}{\small\raise0.25ex\hbox{$\bullet$}}
\setbeamertemplate{enumerate items}[default]
\setbeamerfont{frametitle}{family=\rmfamily,series=\bfseries,size=\LARGE}

\setbeamertemplate{frametitle}{%
  \vspace*{0.55cm}%
  \hspace*{0.15cm}\insertframetitle\par
  \vspace*{0.12cm}%
  \hspace*{0.15cm}{\color{ink}\rule{\dimexpr\paperwidth-1.65cm\relax}{0.6pt}}%
  \vspace*{-0.1cm}%
}
\setbeamertemplate{footline}{%
  \hfill{\color{muted}\scriptsize\insertframenumber\,/\,\inserttotalframenumber}\hspace*{0.75cm}%
  \vspace*{0.35cm}%
}
\setbeamersize{text margin left=0.75cm, text margin right=0.75cm}

\newcommand{\takeaway}[1]{%
  \par\vspace{0.25cm}{\color{accent}\rule[0.2ex]{0.5cm}{2pt}}\hspace{0.25cm}%
  \begin{minipage}[t]{0.9\linewidth}\rmfamily\itshape #1\end{minipage}}
\newcommand{\pv}[1]{\textit{p}\,#1}
\newcommand{\fig}[2][\linewidth]{\includegraphics[width=#1]{../reports/figures/#2.pdf}}
\newcommand{\stat}[2]{{\rmfamily\bfseries\Huge\color{#1}#2}}

\begin{document}

% ---------------------------------------------------------------- 1 title
{
\setbeamercolor{background canvas}{bg=cover}
\setbeamertemplate{footline}{}
\begin{frame}[plain]
  \vspace*{0.6cm}
  {\small\color{muted}International University of Sarajevo $\cdot$ EE418 $\cdot$ Final project, 2026}

  \vspace{1.6cm}
  {\color{accent}\rule{1cm}{2.5pt}}\par\vspace{0.3cm}
  {\rmfamily\bfseries\fontsize{34}{38}\selectfont Detecting Illicit Bitcoin\\[0.1cm]Transactions\par}
  \vspace{0.35cm}
  {\large\color{ink!80}Graphs, imbalance and drift on the Elliptic dataset\par}

  \vspace{1.4cm}
  {\small\textbf{Anes Berbi\'c} (210302141)\hfill
   {\color{muted}github.com/berbicanes/EllipticML}}
\end{frame}
}

% ---------------------------------------------------------------- 2 question + answer
\begin{frame}{The question, and the short answer}
  {\rmfamily\itshape\Large Can we tell illicit Bitcoin transactions apart from licit ones,
  using their features and their position in the transaction graph?\par}
  \vspace{0.6cm}
  \begin{columns}[T]
    \begin{column}{0.32\textwidth}
      \stat{accent}{\MainRFFone}\par\vspace{0.1cm}
      \small illicit F1 on unseen future weeks (Random Forest)
    \end{column}
    \begin{column}{0.32\textwidth}
      \stat{accent}{+0.06}\par\vspace{0.1cm}
      \small F1 from neighbour features: the graph helps, mostly through its simplest summary
    \end{column}
    \begin{column}{0.32\textwidth}
      \stat{illicit}{\DriftRFPost}\par\vspace{0.1cm}
      \small F1 after the dark market shutdown: every model fails
    \end{column}
  \end{columns}
  \takeaway{Detection works, until criminals change behaviour. Most of this talk is
  about measuring that honestly.}
\end{frame}

% ---------------------------------------------------------------- 3 data
\begin{frame}{The data moves under our feet}
  \begin{columns}[c]
    \begin{column}{0.62\textwidth}\fig{eda_labels_over_time}\end{column}
    \begin{column}{0.36\textwidth}\small
      \begin{itemize}\setlength\itemsep{0.35em}
        \item 203,769 transactions, 49 two-week steps, 234,355 flows
        \item Illicit share: \EdaShareTrain{} train, \EdaShareTest{} test, $\approx$0 in steps 43--46
        \item Every step is a separate graph (\EdaComponents{} components)
        \item Illicit neighbours of illicit: \EdaHomIll{} vs \EdaHomLic{} for licit
        \item Illicit transactions never have more than 3 outputs
      \end{itemize}
    \end{column}
  \end{columns}
\end{frame}

% ---------------------------------------------------------------- 4 protocol
\begin{frame}{How we avoided fooling ourselves}
  \small
  \begin{tabular}{p{4.3cm}p{9.3cm}}
    \toprule
    \textbf{Risk} & \textbf{What we did} \\
    \midrule
    Training on the future & Paper's temporal split: steps 1--34 train, 35--49 test \\
    Tuning on the test set & Forward folds inside steps 1--34 pick hyper-parameters \emph{and} threshold \\
    A disguised clock & Time step dropped as a feature (keeping it cut LR's F1 from 0.51 to 0.36) \\
    Leaking through preprocessing & Scaling and SMOTE inside the pipeline; graph features label-free \\
    Lucky seeds & Every model $\times$ 5 seeds \\
    False precision & 95\% CIs by resampling \emph{time steps}, paired tests between models \\
    Silent bugs & Unit tests: no test or future step can reach training \\
    \bottomrule
  \end{tabular}
  \takeaway{First step: reproduce the paper. Our Random Forest: \RepRFAF{} $\pm$ \RepRFAFSd{} vs 0.788 published.}
\end{frame}

% ---------------------------------------------------------------- 5 main comparison
\begin{frame}{Model comparison: two winners, wide error bars}
  \begin{columns}[c]
    \begin{column}{0.58\textwidth}\small
      \begin{tabular}{lccc}
        \toprule
        & F1 & 95\% CI & PR-AUC \\
        \midrule
        Random Forest & \textbf{\MainRFFone} & \MainRFFoneCI & \MainRFPRAUC \\
        XGBoost       & \textbf{\MainXGBFone} & \MainXGBFoneCI & \textbf{\MainXGBPRAUC} \\
        MLP           & \MainMLPFone & \MainMLPFoneCI & \MainMLPPRAUC \\
        Logistic Reg. & \MainLRFone & \MainLRFoneCI & \MainLRPRAUC \\
        \bottomrule
      \end{tabular}
    \end{column}
    \begin{column}{0.4\textwidth}\small
      \begin{itemize}\setlength\itemsep{0.4em}
        \item RF vs XGBoost: paired test \pv{\PairedXGBRFFoneP} $\rightarrow$ \textbf{tied}
        \item Both beat MLP and LR, \pv{<0.001}
        \item An analyst reviewing the top 100 per step finds \MainXGBBudget{} of all illicit
      \end{itemize}
    \end{column}
  \end{columns}
  \takeaway{The interval is 0.27 wide: the test set is 15 time steps, not 16,670
  independent rows. Many published differences on this benchmark are smaller.}
\end{frame}

% ---------------------------------------------------------------- 6 imbalance
\begin{frame}{Imbalance: re-balancing barely matters}
  \fig[0.78\linewidth]{e03_imbalance}
  \takeaway{SMOTE and class weights move PR-AUC by about one point. Threshold tuning
  helps weak models (LR \ImbLRNoneFixedFone{} $\rightarrow$ \ImbLRNoneTunedFone) but not the trees:
  the best threshold drifts over time too.}
\end{frame}

% ---------------------------------------------------------------- 7 graph
\begin{frame}{Does the graph help?}
  \begin{columns}[c]
    \begin{column}{0.56\textwidth}\fig{e04_graph}\end{column}
    \begin{column}{0.42\textwidth}\small
      \begin{itemize}\setlength\itemsep{0.45em}
        \item \textbf{Neighbour features: yes.} Local $\rightarrow$ local + aggregated:
              \GraphRFLFFone{} $\rightarrow$ \GraphRFAFFone{} (\pv{<0.001})
        \item \textbf{More graph: barely.} Structure features and GNN embeddings: $\le$\,0.02 F1, no PR-AUC
        \item \textbf{GNNs alone: no,} but direction matters: our directed DirSAGE
              \GraphDirSAGEgraphFone{} vs GCN \GraphGCNgraphFone{}
      \end{itemize}
    \end{column}
  \end{columns}
\end{frame}

% ---------------------------------------------------------------- 8 hybrid paradox
\begin{frame}{A better GNN makes a worse hybrid}
  \begin{columns}[c]
    \begin{column}{0.46\textwidth}
      \small Weber et al.\ feed GCN embeddings to a Random Forest. We tried the
      embeddings of our stronger DirSAGE instead.\par\vspace{0.4cm}
      \stat{illicit}{$-$\HybRFFoneDrop}\par\vspace{0.1cm}
      \small F1 for the Random Forest (\pv{\HybRFFoneP})
    \end{column}
    \begin{column}{0.48\textwidth}\footnotesize
      \begin{tabular}{@{}lcc@{}}
        \toprule
        & GCN & DirSAGE \\
        \midrule
        PR-AUC, train rows & \MemGCNTrainPRAUC & \textbf{\MemDirSAGETrainPRAUC} \\
        PR-AUC, test rows & \MemGCNTestPRAUC & \MemDirSAGETestPRAUC \\
        RF weight on embeddings & \MemGCNRFImp & \textbf{\MemDirSAGERFImp} \\
        \bottomrule
      \end{tabular}
    \end{column}
  \end{columns}
  \takeaway{The GNN memorised the training labels, the forest copied it. The published
  hybrid works because its GNN is weak; embeddings must be made out-of-fold.}
\end{frame}

% ---------------------------------------------------------------- 9 unsupervised
\begin{frame}{Without labels: crime is typical, not anomalous}
  \begin{columns}[c]
    \begin{column}{0.48\textwidth}\fig{e06_anomaly_scores}\end{column}
    \begin{column}{0.5\textwidth}\fig{e06_label_efficiency}\end{column}
  \end{columns}
  \takeaway{Flag the most anomalous: PR-AUC \IFAFUnsupAnomalousPRAUC{} (worse than random).
  Flag the most \emph{typical}: \IFAFUnsupTypicalPRAUC{}, no labels. XGBoost needs about
  \LabelsToBeatIF{} labels to beat that.}
\end{frame}

% ---------------------------------------------------------------- 10 collapse
\begin{frame}{Every model fails after the dark market shutdown}
  \begin{columns}[c]
    \begin{column}{0.62\textwidth}\fig{e07_f1_per_step}\end{column}
    \begin{column}{0.36\textwidth}\small
      Mean F1 per step, Random Forest:\par\vspace{0.2cm}
      \begin{tabular}{lc}
        steps 35--42 & \textbf{\DriftRFPre} \\
        steps 43--49 & \textbf{\color{illicit}\DriftRFPost} \\
      \end{tabular}\par\vspace{0.4cm}
      The pooled F1 of $\approx$0.8 averages two regimes: near-perfect detection of
      known crime, and almost none of what replaced it.
    \end{column}
  \end{columns}
\end{frame}

% ---------------------------------------------------------------- 11 why
\begin{frame}{Why: the new crime looks licit}
  \centering\fig[0.6\linewidth]{e07_novelty}\par
  \vspace{-0.1cm}\raggedright
  \begin{columns}[T]
    \begin{column}{0.32\textwidth}\small
      \textbf{Closer to licit than to known illicit:} \NovelPre{} before, \textbf{\color{illicit}\NovelPost} after
    \end{column}
    \begin{column}{0.32\textwidth}\small
      \textbf{Features shifted (KS\,$>$\,0.3):} \AdvIllicitKSShare{} for illicit, \AdvLicitKSShare{} for licit
    \end{column}
    \begin{column}{0.32\textwidth}\small
      \textbf{Retraining every step:} post-shutdown F1 \WFStaticPost{} $\rightarrow$ \WFExpandingPost{}
    \end{column}
  \end{columns}
  \takeaway{The bottleneck is labels for the new behaviour, not the model.
  (\ClockCount{} hidden ``clock'' features were removed first.)}
\end{frame}

% ---------------------------------------------------------------- 12 lessons
\begin{frame}{What I would tell a compliance team}
  \begin{enumerate}\setlength\itemsep{0.55em}
    \item Use a tree ensemble on the transaction's own and its neighbours' features.
    \item Report uncertainty over \emph{time}, not over rows; test differences in pairs.
    \item Do not assume crime is anomalous: here it is the most ordinary traffic.
    \item Plan for regime changes: watch the score distribution, label new cases fast.
    \item Check stacked models for leakage: a strong upstream model can poison the downstream one.
  \end{enumerate}
  \vspace{0.3cm}
  {\small\color{muted}Limitations: anonymised features; only 169 illicit test transactions after step 43;
  small GNN search; labels come from heuristics.}
\end{frame}

% ---------------------------------------------------------------- 13 thanks
{
\setbeamercolor{background canvas}{bg=cover}
\begin{frame}
  \vspace*{0.4cm}
  {\color{accent}\rule{1cm}{2.5pt}}\par\vspace{0.3cm}
  {\rmfamily\bfseries\fontsize{34}{38}\selectfont Thank you\par}
  \vspace{0.3cm}
  {\large Questions are welcome.\par}
  \vspace{1.2cm}
  \small
  \textbf{Code, results, report:} github.com/berbicanes/EllipticML\\[0.15cm]
  {\color{muted}\texttt{make all} rebuilds every number from the raw data;
  every figure in this deck is generated by the pipeline.}\par
  \vspace{0.9cm}
  {\scriptsize\color{muted}
  Weber et al. (2019), Anti-Money Laundering in Bitcoin, KDD Workshop, arXiv:1908.02591 $\cdot$
  Kipf \& Welling (2017), ICLR $\cdot$ Hamilton et al. (2017), NeurIPS $\cdot$
  Rossi et al. (2023), LoG $\cdot$ Liu et al. (2008), ICDM $\cdot$ Lundberg \& Lee (2017), NeurIPS\par}
\end{frame}
}

\end{document}
